\section{Experimental Setup and Results}
\label{sec:experiments}

\subsection{Experimental setup}
\label{sec:setup}

\textbf{Pre-training data and splits.}
Pre-training uses Tiny-ImageNet~\cite{le2015tinyimagenet}: 100\,000 training
images, 200 classes, native $64\times64$ resolution
(Figure~\ref{fig:datasets}); labels are never used during pre-training. The
pre-training set sizes are $N \in \{1\mathrm{k}, 2\mathrm{k}, 4\mathrm{k},
8\mathrm{k}, 16\mathrm{k}, 32\mathrm{k}, 64\mathrm{k}, 100\mathrm{k}\}$,
drawn as prefixes of a single shuffled master list (seed 42) that interleaves
one image per class, so every split is exactly class-balanced and smaller
splits are subsets of larger ones. The nesting means a larger split differs
from a smaller one only by the images added, so curve differences are
attributable to $N$ alone.

\textbf{Budgets and validation probe.}
Each split trains a fixed epoch budget of \{1k--8k: 1000, 16k: 600, 32k: 400,
64k: 300, 100k: 250\} epochs. Budgets shrink with $N$ because small splits
reach their validation peak only after hundreds of epochs while large splits
peak within tens; each budget is sized to be generously past the peak, not to
a convergence criterion, and runs are extendable when a split's validation is
still rising at its budget end (disclosed per split in
Section~\ref{sec:results_curve}). The validation probe runs every 4 epochs:
the Tiny-ImageNet \texttt{valid} split (10\,000 labelled images held out from
all pre-training splits) is partitioned class-stratified into a 5\,000-image
gallery and 5\,000-image query set (25 per class each), and the query images
are classified by their $k{=}20$ nearest gallery neighbours (Euclidean
distance, unweighted) in CLS-embedding space. The same probe also records the
effective rank of the query embeddings, as the dimensional-collapse
diagnostic of Section~\ref{sec:results_overtraining} --- logged only, never
used for selection.

\textbf{Downstream evaluation.}
The annealed encoder is frozen, and a single linear layer
($192\rightarrow10$) is trained on CIFAR-10~\cite{krizhevsky2009cifar},
resized $32\rightarrow64$ (bilinear) and ImageNet-normalized to match
pre-training statistics. The probe protocol is shared across the five-student
project: SGD with momentum 0.9, no weight decay, batch 512, learning rate
$0.2$ ($0.1 \times \mathrm{batch}/256$, the linear scaling rule), cosine
annealed to zero over 100 epochs, fixed evaluation seed independent of the
pre-training seed. Two protocol details are disclosed: the training
transform includes a random horizontal flip, and we report the best test
top-1 accuracy across the 100 probe epochs --- a metric that peeks at the
test curve; the final-epoch accuracy is recorded alongside and the gap
reported (Section~\ref{sec:limitations}).

\textbf{Baselines.}
A random-initialization floor applies the identical probe protocol to an
untrained encoder; it isolates how much of the linear-probe accuracy is
bought by pre-training rather than by the architecture and the probe itself.

\textbf{Seeds, tuning, and provenance.}
Every (split, seed) pre-training run is indexed by the code commit, the split
size, and the seed; the campaign uses seeds 42, 43, 44, and curves report
mean and spread over seeds. One hyperparameter was tuned: the learning rate,
by a one-off ablation on the 8k split over $\{1, 3, 5, 7, 15\} \times
10^{-4}$. The result is an inverted U: $10^{-4}$ under-trains, $3$--$7
\times 10^{-4}$ form a plateau within probe noise, and $1.5 \times 10^{-3}$
degrades both accuracy and effective rank. We picked $5 \times 10^{-4}$ ---
the plateau center, and the reference value of DINO at batch
256~\cite{caron2021dino} --- and use this single rate for all splits, keeping
$N$ the only variable that changes across the curve. All remaining
hyperparameters are fixed by the shared protocol or by the published recipes
cited in Section~\ref{sec:method}. All runs execute on one consumer GPU;
wall-clock costs are reported in Section~\ref{sec:discussion}.

\begin{figure}[t]
  \centering
  % [PLACEHOLDER --- replaced by figures/dataset_samples.pdf in the figure pass]
  \fbox{\parbox[c][3.2cm][c]{0.92\linewidth}{\centering\small
    Dataset samples (pending): a grid of Tiny-ImageNet images ($64\times64$,
    200 classes) next to CIFAR-10 images ($32\times32$, 10 classes, shown
    resized to $64\times64$).}}
  \caption{Pre-training and downstream data. Tiny-ImageNet (left) provides
  the unlabelled pre-training images at $64\times64$; CIFAR-10 (right),
  resized to the same resolution, is the downstream classification task. The
  two datasets share no images, so the linear probe measures transfer of the
  learned representation, not recall of the pre-training set.}
  \label{fig:datasets}
\end{figure}

% ---------------------------------------------------------------------------
% [ALL PENDING] Question-headed result subsections (WE1/WE2); written as
% hypothesis-style stubs until the campaign finishes; interpretations flagged.
% ---------------------------------------------------------------------------

\subsection{How does downstream accuracy scale with pre-training set size?}
\label{sec:results_curve}
[PENDING --- data-efficiency curve over \{1k...100k\} with random-init floor
(and supervised ceiling if run); per split, disclose whether the validation
signal plateaued within the budget.]

\subsection{At the smallest scales, does Barlow Twins over-train --- and is it collapse?}
\label{sec:results_overtraining}
[PENDING --- long-run accuracy and effective-rank traces on the smallest
split(s); observation stated firmly, the overfitting mechanism labelled as an
interpretation we did not isolate.]

\subsection{Does projector width interact with dataset size?}
\label{sec:results_projector}
[PENDING --- accuracy vs projector dimension \{256, 512, 1024, 2048\} on the
smaller splits, $\geq 2$ seeds; report a shift of the best width with $N$, or
the null result.]

\subsection{Does the in-domain selection signal track downstream accuracy?}
\label{sec:results_correlation}
[PENDING, optional --- kNN(Tiny-ImageNet) vs CIFAR-10 probe correlation across
checkpoints from the diagnostic pass.]
